{\color{red}
\section{Open questions for the author (NOTE list)}
\noindent[\textbf{New section.} Points that the fix notes do not determine uniquely, or that are not settled; in each case the text of the file was left as the author has it (or as the fix note gives it, where the fix note is unambiguous), with a red \textbf{NOTE} at the place.]
\begin{enumerate}
\item \textbf{DGLAP subtraction term} (notes at Group II Row IV and in the definitions at the beginning of the finite terms; the DGLAP term itself is in \texttt{Final\_terms\_fixed}). fix\_poles\_dglap Step 11 writes $-\frac{2N_{c}}{\xi}\mathcal{F}(0)$; this equals the author's subtraction $-2N_{c}\frac{1-\xi}{\xi}\mathcal{F}(0)$ \emph{plus} the remnant of Group II Row IV, while Step 11 also keeps the remnant in Row IV. The author's convention (remnant in Row IV) is kept. Which writing is wanted for the paper?
\item \textbf{Fragmentation-function scheme} (\texttt{Final\_terms\_fixed}). \textbf{Answered} (author's answer (date 2026-11-20): ``Yes''.) The fragmentation functions are $\overline{\rm MS}$ number densities (pure-pole renormalisation with the $d$-dimensional Jacobian $\xi^{-d}$), i.e.\ the author's $\log\frac{(y'-y)^{2}\mu^{2}}{4\xi^{2}}$ form; this is now stated as the definition, and $\Delta$ to the CXY form is kept as a remark. (Not done, a cross-check only: an independent comparison with a standard NLO collinear $pp\to gX$ calculation.)
\item \textbf{Coordinate or local form of DGLAP and RC} (\texttt{Final\_terms\_fixed}). \textbf{Answered} (author's answer (date 2026-11-20): ``write it in both forms''.) Both forms are given, with their equivalence (DGLAP and RC sections of \texttt{Final\_terms\_fixed}): the coordinate forms are separately logarithmically IR divergent at large parent position, their sum with Group II Row I (second part) and Group III Row III is not (the tail cancels term by term, \texttt{notes/fix\_tail.tex}); in the local form the remainder $N_{\mathcal{R}}$ carries the tail.
\item \textbf{Group II Row IV large-distance tail} (Group II Rows I, IV; Group III Row III). \textbf{Settled} in \texttt{notes/fix\_tail.tex}: the term-by-term check on the finite terms as written is done; the tails of Group II Row IV, Group II Row I (second part) and Group III Row III (second part, with c.c.) are $-\frac{11}{6}+\frac16\cos2\varphi$, $-\frac{11}{6}+\frac16\cos2\varphi$, $+\frac{11}{3}-\frac13\cos2\varphi$ (units $2N_{c}\mathcal{C}\frac{4\pi}{k^{2}}\int\frac{d^{2}Y}{Y^{2}}G_{\infty}$) and cancel, also in the BFKL and JIMWLK sectors; the three rows must be integrated together over the pair position. [Consistency update: this item still read ``pending'' after the fix\_tail update.]
\item \textbf{Prescription of the energy denominators of $\mf{B}_{3}$} (Group I Row III; block $\mathcal{S}$). \textbf{Answered} (author's answer (date 2026-11-20): ``yes''.) The energy denominators are principal values; this is stated as the prescription, the $\pm i\epsilon$ alternative is not used.
\item \textbf{Order of limits} (Finite terms, intro). The remaining pole $-\frac{\ell}{2}\frac{2}{\epsilon}$, $\ell=\log(1-\frac{k^{+}}{\vee})$, vanishes only if $\vee\to\infty$ is taken before $\epsilon\to0$; this is assumed and should be stated in the paper.
\item \textbf{Cubic contact terms} (singlet section). The cubic rows leave quadratic contact terms after the projectile average; whether they cancel among the rows, or against the contact terms of the three-$\rho$ evolution terms, was not checked.
\item \textbf{Integrability at the contact configurations} (singlet section): checked only for two kernels.
\item \textbf{Assumption (A2)} (singlet section): no C-odd ($d^{abc}$) or chiral ($f^{abc}g_{f}$) three-point function of the projectile; otherwise the symmetrised cubic parts, odd in $k$, are added.
\item \textbf{Kinematically constrained reading} of the soft block (fix\_soft Step 7). \textbf{Answered} (author's answer (date 2026-11-20): ``write it in both form''.) The kinematically constrained form is written out in full, Eqs.~\eqref{eq:fx-FLam}, \eqref{eq:fx-Skc}, with its equivalence to the paired form, Eq.~\eqref{eq:fx-kcequiv}. Remaining \textbf{NOTE}: for pair D the identification of $\log\frac{k^{+}}{\Lambda}\alpha_{D}V_{D}C_{D}[\kappa_{D}(w)-\kappa_{D}(w')]$ with JIMWLK-list terms is taken from fix\_soft Step 7, not re-derived (the Group I member of pair D has no $\log\frac{k^{+}}{\Lambda}$ term of its own).
\item \textbf{Not done in fix\_soft}: complete numerical integration with Wilson lines; re-check of the regular terms (taken from \texttt{Finite\_terms\_final}); the non-singlet $\log\frac{k^{+}}{\Lambda}\int\frac{d^{2}z}{z^{2}}C_{A}$ of the JIMWLK formula.
\end{enumerate}}\par
